\chapter*{INTRODUCTION}
\addcontentsline{toc}{chapter}{INTRODUCTION}
\label{ch:intro}

Industrial control systems (ICS) are the interface between digital policy and physical plant: programmable logic controllers, supervisory control and data acquisition (SCADA) servers, and remote terminal units that translate operator intent into actuator commands and report sensor readings back.  Unlike enterprise IT, where a successful intrusion typically results in data theft or service disruption, an undetected intrusion into an ICS can cause damage in the physical world that software-only compromises cannot --- damage to equipment, environmental release, or harm to human operators.  The historical record establishes that this is not a hypothetical concern.  Stuxnet~\cite{langner2013stuxnet} demonstrated that targeted control-loop manipulation can destroy industrial centrifuges while presenting normal-looking telemetry to monitoring stations.  The 2015 Ukraine power-grid attack~\cite{ukraine2016analysis} caused a multi-hour outage affecting roughly 225{,}000 customers across three regional distribution operators, with attackers using legitimate SCADA functionality to open breakers from inside the operators' own control consoles.  Earlier incidents such as the Maroochy Shire wastewater spill (2000) showed that even comparatively unsophisticated insider-style attacks can cause significant environmental harm when basic anomaly-detection telemetry is absent.

The defensive response from the research community has been to construct labelled testbeds and publish anomaly detectors evaluated on them.  Public ICS datasets --- HAI~\cite{shin2020hai}, SWaT~\cite{goh2017swat}, WADI~\cite{ahmed2017wadi}, the Morris gas-pipeline corpus~\cite{morris2014industrial} --- enable reproducible benchmarking, and a large body of deep-learning anomaly detectors reports high F\textsubscript{1} scores on them.  Recent transformer- and adversarial-autoencoder-based detectors (TranAD~\cite{tuli2022tranad}, USAD~\cite{audibert2020usad}) report point-adjusted F\textsubscript{1} above $0.95$ on HAI and SWaT, which would suggest a near-solved problem.  An operator considering deployment, however, faces three questions that the headline numbers skip over:

\begin{enumerate}
  \item \textbf{Do the rankings survive time-aware evaluation?}  Point-wise F\textsubscript{1} treats every timestep as independent.  Point-adjust F\textsubscript{1} (PA-F\textsubscript{1})~\cite{xu2018unsupervised} inflates scores by granting full credit when any single point inside an attack event is flagged --- a recent critique~\cite{kim2022towards} shows that under PA-F\textsubscript{1} a random detector can outperform genuine ones.  eTaPR~\cite{hwang2022etapr} attempts to restore event-level rigour by weighting detection coverage and timeliness.  Which metric one reports changes which model ``wins.''
  \item \textbf{Do these detectors generalize across testbeds?}  All cited papers train and test on the same dataset.  An HAI-trained model should, in principle, capture ICS-general structure and transfer to an industrial gas-pipeline testbed --- or the claim of ``anomaly detection'' is overstated.  A detector that only works on the testbed it was trained on is, in operational terms, a one-off model fitted to a particular plant.
  \item \textbf{Can the detector localize?}  A flagged window is actionable only if an operator can see \emph{which sensor} is implicated.  Detection without localization forces a manual triage step that scales poorly when alerts fire across hundreds of sensors.  We evaluate per-feature attribution against HAI's per-process attack labels and ask whether the best-detecting model is also the best-localizing one.
\end{enumerate}

These three questions motivate the diploma project.  The formal statement of the task --- purpose, objectives, object and subject of research, scientific novelty and practical significance, and structure of the work --- is given as the closing section of Chapter~\ref{ch:background}, after the analytical survey of methods, datasets, and metrics that frames it (Section~\ref{sec:task-statement}).
